% ======================================================================
% Kap. 6 – Abbildungen (AP1). Entwurf für die Zusammenführung in Kapitel/06_Konzept.tex (AP8).
% Jede Abbildung wird im Fließtext VOR dem Float mit \ref eingeführt und danach erläutert
% (CLAUDE.md "Nie ohne Textbezug"); die figure-Umgebung steht direkt nach dem Absatz mit dem ersten Verweis.
% Alle Grafiken: eigene Darstellung (keine Quellenangabe in der Caption), Stil ba-grafik.tex, nicht skaliert.
% Die Datenbasis jeder Grafik steht als Kommentar am Anfang ihrer Datei in Bilder/tikz/.
% ======================================================================

% ---- 6.2.1 Komponenten und Datenflüsse (Kern) -------------------------------------------------
% Text davor: Abbildung~\ref{fig:architektur}, von oben nach unten durchgehen.
\begin{figure}[htbp]
  \centering
  % Gesamtarchitektur der Umgebung: Bedienwege, Funktionsknoten, Herstellerschnittstelle, Simulation oder Roboter.
% Eigene Darstellung (Art E, Kap. 6.2.1). Keine Quelle in der Caption.
% Basis: G1@25d5c1a g1pilot/docs/02-T-Architektur.md (Gesamtbild, Laufzeit-Topologie, Zwei Kommunikationsebenen,
%   Sim <-> Real: Ein Schalter, Joint-Layout/Zuständigkeit der Gelenke), docs/05-T-Nodes-und-Topics.md,
%   docs/10-T-Demo-GUI.md (Grundsatz), docs/30-T-Stehen-und-Gehen.md (Warum zwei Implementierungen),
%   docs/21-T-Arm-API.md (Arm-API als eigener Knoten), docs/50-A-Echter-Roboter.md Abschnitt E und
%   docs/22-T-Inspire-Haende.md (Hände real getrennter Anschluss, Sim über Simulation).
% Vereinfachungen: Navigation (nur Simulation) und Hand-Profile nicht gezeichnet; die Bedienoberfläche läuft
%   technisch im selben Container wie die Funktionsknoten, steht hier als Bedienweg oben.
% Stil: ba-grafik.tex. Einbinden: % Gesamtarchitektur der Umgebung: Bedienwege, Funktionsknoten, Herstellerschnittstelle, Simulation oder Roboter.
% Eigene Darstellung (Art E, Kap. 6.2.1). Keine Quelle in der Caption.
% Basis: G1@25d5c1a g1pilot/docs/02-T-Architektur.md (Gesamtbild, Laufzeit-Topologie, Zwei Kommunikationsebenen,
%   Sim <-> Real: Ein Schalter, Joint-Layout/Zuständigkeit der Gelenke), docs/05-T-Nodes-und-Topics.md,
%   docs/10-T-Demo-GUI.md (Grundsatz), docs/30-T-Stehen-und-Gehen.md (Warum zwei Implementierungen),
%   docs/21-T-Arm-API.md (Arm-API als eigener Knoten), docs/50-A-Echter-Roboter.md Abschnitt E und
%   docs/22-T-Inspire-Haende.md (Hände real getrennter Anschluss, Sim über Simulation).
% Vereinfachungen: Navigation (nur Simulation) und Hand-Profile nicht gezeichnet; die Bedienoberfläche läuft
%   technisch im selben Container wie die Funktionsknoten, steht hier als Bedienweg oben.
% Stil: ba-grafik.tex. Einbinden: % Gesamtarchitektur der Umgebung: Bedienwege, Funktionsknoten, Herstellerschnittstelle, Simulation oder Roboter.
% Eigene Darstellung (Art E, Kap. 6.2.1). Keine Quelle in der Caption.
% Basis: G1@25d5c1a g1pilot/docs/02-T-Architektur.md (Gesamtbild, Laufzeit-Topologie, Zwei Kommunikationsebenen,
%   Sim <-> Real: Ein Schalter, Joint-Layout/Zuständigkeit der Gelenke), docs/05-T-Nodes-und-Topics.md,
%   docs/10-T-Demo-GUI.md (Grundsatz), docs/30-T-Stehen-und-Gehen.md (Warum zwei Implementierungen),
%   docs/21-T-Arm-API.md (Arm-API als eigener Knoten), docs/50-A-Echter-Roboter.md Abschnitt E und
%   docs/22-T-Inspire-Haende.md (Hände real getrennter Anschluss, Sim über Simulation).
% Vereinfachungen: Navigation (nur Simulation) und Hand-Profile nicht gezeichnet; die Bedienoberfläche läuft
%   technisch im selben Container wie die Funktionsknoten, steht hier als Bedienweg oben.
% Stil: ba-grafik.tex. Einbinden: \input{Bilder/tikz/architektur-gesamt} in einer figure-Umgebung. Nicht skalieren.
\begin{tikzpicture}[ba]
  \tikzset{
    bw/.style={ba block, text width=31mm, minimum height=10mm},
    fk/.style={ba block hervor, minimum height=17mm, minimum width=0pt},
    sub/.style={ba block, font=\footnotesize, minimum height=10mm, minimum width=0pt, inner sep=2pt},
    ziel/.style={ba pfeil},
    zust/.style={ba pfeil, dash pattern=on 3pt off 2pt},
    klein/.style={font=\footnotesize, text=baGrau},
  }
  % ---------------- Ebene 1: Bedienwege ----------------
  \node[klein, anchor=south west] at (1.3,11.45) {Bedienwege (rechnen keine Bewegung, senden Ziele und lesen Zustände)};
  \node[bw] (B1) at (3.0,10.65) {Bedien-\\oberfläche};
  \node[bw] (B2) at (6.6,10.65) {Gamepad};
  \node[bw] (B3) at (10.2,10.65) {Smartphone};
  \node[bw] (B4) at (13.8,10.65) {Externe\\Programme};
  % Sammelleitung ROS 2
  \draw[ba normal] (3.0,9.95) -- (13.8,9.95);
  \foreach \x in {3.0,6.6,10.2,13.8} {\draw[ba normal] (\x,10.15) -- (\x,9.95);}
  \node[klein, fill=white, inner sep=1.5pt] at (4.3,9.95) {ROS 2};
  % ---------------- Container Umgebung ----------------
  \draw[ba hilfslinie] (1.2,5.25) rectangle (15.6,9.35);
  \node[klein, anchor=south west] at (1.3,5.3) {Container Umgebung};
  % Ebene 2: Funktionsknoten
  \node[fk, text width=22mm] (RZ) at (2.5,7.4) {Roboter-\\zustand};
  \node[fk, text width=28mm] (AR) at (5.9,7.4) {Armregelung\\mit Planung};
  \node[fk, minimum width=51mm, minimum height=24mm, anchor=center] (FB) at (10.55,7.4) {};
  \node[font=\small, anchor=north] at (10.55,8.55) {Fortbewegung};
  \node[sub, text width=21mm] (ER) at (9.2,7.15) {Ersatzregler\\(Simulation)};
  \node[sub, text width=21mm] (OA) at (11.9,7.15) {Anbindung Onboard-Regler (Realbetrieb)};
  \node[fk, text width=19mm] (HA) at (14.35,7.4) {Hände};
  % Pfeile Bedienwege -> Funktionsknoten (Ziele nach unten, Zustände nach oben)
  \draw[ziel] (5.6,9.95) -- (5.6,8.6);
  \draw[zust] (6.2,8.6) -- (6.2,9.95);
  \draw[ziel] (8.6,9.95) -- (8.6,8.6);
  \draw[zust] (12.5,8.6) -- (12.5,9.95);
  \draw[ziel] (14.05,9.95) -- (14.05,8.6);
  \draw[zust] (14.65,8.6) -- (14.65,9.95);
  \draw[zust] (2.5,8.25) -- (2.5,9.95);
  % Schutzmechanismen
  \node[ba block hervor, minimum height=6mm, minimum width=108mm, font=\footnotesize] (SM) at (10.1,5.95)
    {Schutzmechanismen: begrenzen \textperiodcentered{} vermeiden \textperiodcentered{} überwachen \textperiodcentered{} anhalten};
  \draw[ba duenn, color=baGruen] (5.9,6.55) -- (5.9,6.25);
  \draw[ba duenn, color=baGruen] (10.55,6.55) -- (10.55,6.25);
  \draw[ba duenn, color=baGruen] (14.35,6.55) -- (14.35,6.25);
  % ---------------- Ebene 3: Herstellerschnittstelle ----------------
  \node[ba block, fill=baHellgrau!40, minimum height=8mm, minimum width=143mm] (DDS) at (8.4,4.2)
    {Herstellerschnittstelle (Unitree-DDS)};
  % Befehle nach unten, Zustände nach oben zwischen Funktionsknoten und Schnittstelle
  \draw[zust] (2.5,4.6) -- (2.5,6.55);
  \draw[ziel] (5.6,5.65) -- (5.6,4.6);
  \draw[zust] (6.2,4.6) -- (6.2,5.65);
  \draw[ziel] (9.5,5.65) -- (9.5,4.6);
  \draw[zust] (10.1,4.6) -- (10.1,5.65);
  \draw[ziel] (11.5,5.65) -- (11.5,4.6);
  \draw[zust] (12.1,4.6) -- (12.1,5.65);
  \draw[ziel] (14.05,5.65) -- (14.05,4.6);
  \draw[zust] (14.65,4.6) -- (14.65,5.65);
  % ---------------- Umschalter ----------------
  \node[draw=baAkzent, ba normal, fill=white, rounded corners=1.5pt, text=black, align=center,
        minimum height=9mm, inner sep=3pt] (SW) at (8.4,2.75) {Umschalter\\Simulation / Realbetrieb};
  \draw[baAkzent, ba normal] (8.4,3.8) -- (SW.north);
  % ---------------- Ebene 4: Simulation oder Roboter ----------------
  \draw[ba gruppe] (1.3,-1.1) rectangle (7.0,1.5);
  \node[klein, anchor=north west] at (1.4,1.5) {Container Simulation};
  \node[ba block, minimum width=0pt, text width=22mm] (SIM) at (2.9,-0.1) {MuJoCo-Simulation};
  \node[ba block, minimum width=0pt, text width=22mm] (SH) at (5.5,-0.1) {Simulierte Hände};
  \draw[ba gruppe] (9.8,-1.1) rectangle (15.5,1.5);
  \node[klein, anchor=north west] at (9.9,1.5) {Realbetrieb};
  \node[ba block, minimum width=0pt, text width=22mm] (RB) at (11.4,-0.1) {Unitree G1 mit Onboard-Regler};
  \node[ba block, minimum width=0pt, text width=22mm] (RH) at (14.0,-0.1) {Inspire-Hände};
  \draw[baAkzent, ba normal, <->] (SW.west) -| (4.15,1.5);
  \draw[baAkzent, ba normal, <->] (SW.east) -| (12.65,1.5);
  \node[font=\footnotesize, anchor=east, align=right] at (4.05,2.1) {Domain 1,\\Schnittstelle Loopback};
  \node[font=\footnotesize, anchor=west, align=left] at (12.75,2.1) {Domain 0,\\Netzwerkkarte};
  % Hände real: eigener Anschluss neben der Herstellerschnittstelle
  \draw[ba normal, dash pattern=on 1pt off 2pt, <->] (HA.east) -- (15.85,7.4) -- (15.85,-0.1) -- (RH.east);
  \node[klein, anchor=east, align=right, fill=white, inner sep=1pt] at (15.8,3.2) {eigener Anschluss,\\nur Realbetrieb};
  % Legende
  \draw[ziel] (1.3,-1.75) -- (2.2,-1.75);
  \node[klein, anchor=west] at (2.3,-1.75) {Ziele und Befehle};
  \draw[zust] (6.0,-1.75) -- (6.9,-1.75);
  \node[klein, anchor=west] at (7.0,-1.75) {Zustände und Meldungen};
\end{tikzpicture}
 in einer figure-Umgebung. Nicht skalieren.
\begin{tikzpicture}[ba]
  \tikzset{
    bw/.style={ba block, text width=31mm, minimum height=10mm},
    fk/.style={ba block hervor, minimum height=17mm, minimum width=0pt},
    sub/.style={ba block, font=\footnotesize, minimum height=10mm, minimum width=0pt, inner sep=2pt},
    ziel/.style={ba pfeil},
    zust/.style={ba pfeil, dash pattern=on 3pt off 2pt},
    klein/.style={font=\footnotesize, text=baGrau},
  }
  % ---------------- Ebene 1: Bedienwege ----------------
  \node[klein, anchor=south west] at (1.3,11.45) {Bedienwege (rechnen keine Bewegung, senden Ziele und lesen Zustände)};
  \node[bw] (B1) at (3.0,10.65) {Bedien-\\oberfläche};
  \node[bw] (B2) at (6.6,10.65) {Gamepad};
  \node[bw] (B3) at (10.2,10.65) {Smartphone};
  \node[bw] (B4) at (13.8,10.65) {Externe\\Programme};
  % Sammelleitung ROS 2
  \draw[ba normal] (3.0,9.95) -- (13.8,9.95);
  \foreach \x in {3.0,6.6,10.2,13.8} {\draw[ba normal] (\x,10.15) -- (\x,9.95);}
  \node[klein, fill=white, inner sep=1.5pt] at (4.3,9.95) {ROS 2};
  % ---------------- Container Umgebung ----------------
  \draw[ba hilfslinie] (1.2,5.25) rectangle (15.6,9.35);
  \node[klein, anchor=south west] at (1.3,5.3) {Container Umgebung};
  % Ebene 2: Funktionsknoten
  \node[fk, text width=22mm] (RZ) at (2.5,7.4) {Roboter-\\zustand};
  \node[fk, text width=28mm] (AR) at (5.9,7.4) {Armregelung\\mit Planung};
  \node[fk, minimum width=51mm, minimum height=24mm, anchor=center] (FB) at (10.55,7.4) {};
  \node[font=\small, anchor=north] at (10.55,8.55) {Fortbewegung};
  \node[sub, text width=21mm] (ER) at (9.2,7.15) {Ersatzregler\\(Simulation)};
  \node[sub, text width=21mm] (OA) at (11.9,7.15) {Anbindung Onboard-Regler (Realbetrieb)};
  \node[fk, text width=19mm] (HA) at (14.35,7.4) {Hände};
  % Pfeile Bedienwege -> Funktionsknoten (Ziele nach unten, Zustände nach oben)
  \draw[ziel] (5.6,9.95) -- (5.6,8.6);
  \draw[zust] (6.2,8.6) -- (6.2,9.95);
  \draw[ziel] (8.6,9.95) -- (8.6,8.6);
  \draw[zust] (12.5,8.6) -- (12.5,9.95);
  \draw[ziel] (14.05,9.95) -- (14.05,8.6);
  \draw[zust] (14.65,8.6) -- (14.65,9.95);
  \draw[zust] (2.5,8.25) -- (2.5,9.95);
  % Schutzmechanismen
  \node[ba block hervor, minimum height=6mm, minimum width=108mm, font=\footnotesize] (SM) at (10.1,5.95)
    {Schutzmechanismen: begrenzen \textperiodcentered{} vermeiden \textperiodcentered{} überwachen \textperiodcentered{} anhalten};
  \draw[ba duenn, color=baGruen] (5.9,6.55) -- (5.9,6.25);
  \draw[ba duenn, color=baGruen] (10.55,6.55) -- (10.55,6.25);
  \draw[ba duenn, color=baGruen] (14.35,6.55) -- (14.35,6.25);
  % ---------------- Ebene 3: Herstellerschnittstelle ----------------
  \node[ba block, fill=baHellgrau!40, minimum height=8mm, minimum width=143mm] (DDS) at (8.4,4.2)
    {Herstellerschnittstelle (Unitree-DDS)};
  % Befehle nach unten, Zustände nach oben zwischen Funktionsknoten und Schnittstelle
  \draw[zust] (2.5,4.6) -- (2.5,6.55);
  \draw[ziel] (5.6,5.65) -- (5.6,4.6);
  \draw[zust] (6.2,4.6) -- (6.2,5.65);
  \draw[ziel] (9.5,5.65) -- (9.5,4.6);
  \draw[zust] (10.1,4.6) -- (10.1,5.65);
  \draw[ziel] (11.5,5.65) -- (11.5,4.6);
  \draw[zust] (12.1,4.6) -- (12.1,5.65);
  \draw[ziel] (14.05,5.65) -- (14.05,4.6);
  \draw[zust] (14.65,4.6) -- (14.65,5.65);
  % ---------------- Umschalter ----------------
  \node[draw=baAkzent, ba normal, fill=white, rounded corners=1.5pt, text=black, align=center,
        minimum height=9mm, inner sep=3pt] (SW) at (8.4,2.75) {Umschalter\\Simulation / Realbetrieb};
  \draw[baAkzent, ba normal] (8.4,3.8) -- (SW.north);
  % ---------------- Ebene 4: Simulation oder Roboter ----------------
  \draw[ba gruppe] (1.3,-1.1) rectangle (7.0,1.5);
  \node[klein, anchor=north west] at (1.4,1.5) {Container Simulation};
  \node[ba block, minimum width=0pt, text width=22mm] (SIM) at (2.9,-0.1) {MuJoCo-Simulation};
  \node[ba block, minimum width=0pt, text width=22mm] (SH) at (5.5,-0.1) {Simulierte Hände};
  \draw[ba gruppe] (9.8,-1.1) rectangle (15.5,1.5);
  \node[klein, anchor=north west] at (9.9,1.5) {Realbetrieb};
  \node[ba block, minimum width=0pt, text width=22mm] (RB) at (11.4,-0.1) {Unitree G1 mit Onboard-Regler};
  \node[ba block, minimum width=0pt, text width=22mm] (RH) at (14.0,-0.1) {Inspire-Hände};
  \draw[baAkzent, ba normal, <->] (SW.west) -| (4.15,1.5);
  \draw[baAkzent, ba normal, <->] (SW.east) -| (12.65,1.5);
  \node[font=\footnotesize, anchor=east, align=right] at (4.05,2.1) {Domain 1,\\Schnittstelle Loopback};
  \node[font=\footnotesize, anchor=west, align=left] at (12.75,2.1) {Domain 0,\\Netzwerkkarte};
  % Hände real: eigener Anschluss neben der Herstellerschnittstelle
  \draw[ba normal, dash pattern=on 1pt off 2pt, <->] (HA.east) -- (15.85,7.4) -- (15.85,-0.1) -- (RH.east);
  \node[klein, anchor=east, align=right, fill=white, inner sep=1pt] at (15.8,3.2) {eigener Anschluss,\\nur Realbetrieb};
  % Legende
  \draw[ziel] (1.3,-1.75) -- (2.2,-1.75);
  \node[klein, anchor=west] at (2.3,-1.75) {Ziele und Befehle};
  \draw[zust] (6.0,-1.75) -- (6.9,-1.75);
  \node[klein, anchor=west] at (7.0,-1.75) {Zustände und Meldungen};
\end{tikzpicture}
 in einer figure-Umgebung. Nicht skalieren.
\begin{tikzpicture}[ba]
  \tikzset{
    bw/.style={ba block, text width=31mm, minimum height=10mm},
    fk/.style={ba block hervor, minimum height=17mm, minimum width=0pt},
    sub/.style={ba block, font=\footnotesize, minimum height=10mm, minimum width=0pt, inner sep=2pt},
    ziel/.style={ba pfeil},
    zust/.style={ba pfeil, dash pattern=on 3pt off 2pt},
    klein/.style={font=\footnotesize, text=baGrau},
  }
  % ---------------- Ebene 1: Bedienwege ----------------
  \node[klein, anchor=south west] at (1.3,11.45) {Bedienwege (rechnen keine Bewegung, senden Ziele und lesen Zustände)};
  \node[bw] (B1) at (3.0,10.65) {Bedien-\\oberfläche};
  \node[bw] (B2) at (6.6,10.65) {Gamepad};
  \node[bw] (B3) at (10.2,10.65) {Smartphone};
  \node[bw] (B4) at (13.8,10.65) {Externe\\Programme};
  % Sammelleitung ROS 2
  \draw[ba normal] (3.0,9.95) -- (13.8,9.95);
  \foreach \x in {3.0,6.6,10.2,13.8} {\draw[ba normal] (\x,10.15) -- (\x,9.95);}
  \node[klein, fill=white, inner sep=1.5pt] at (4.3,9.95) {ROS 2};
  % ---------------- Container Umgebung ----------------
  \draw[ba hilfslinie] (1.2,5.25) rectangle (15.6,9.35);
  \node[klein, anchor=south west] at (1.3,5.3) {Container Umgebung};
  % Ebene 2: Funktionsknoten
  \node[fk, text width=22mm] (RZ) at (2.5,7.4) {Roboter-\\zustand};
  \node[fk, text width=28mm] (AR) at (5.9,7.4) {Armregelung\\mit Planung};
  \node[fk, minimum width=51mm, minimum height=24mm, anchor=center] (FB) at (10.55,7.4) {};
  \node[font=\small, anchor=north] at (10.55,8.55) {Fortbewegung};
  \node[sub, text width=21mm] (ER) at (9.2,7.15) {Ersatzregler\\(Simulation)};
  \node[sub, text width=21mm] (OA) at (11.9,7.15) {Anbindung Onboard-Regler (Realbetrieb)};
  \node[fk, text width=19mm] (HA) at (14.35,7.4) {Hände};
  % Pfeile Bedienwege -> Funktionsknoten (Ziele nach unten, Zustände nach oben)
  \draw[ziel] (5.6,9.95) -- (5.6,8.6);
  \draw[zust] (6.2,8.6) -- (6.2,9.95);
  \draw[ziel] (8.6,9.95) -- (8.6,8.6);
  \draw[zust] (12.5,8.6) -- (12.5,9.95);
  \draw[ziel] (14.05,9.95) -- (14.05,8.6);
  \draw[zust] (14.65,8.6) -- (14.65,9.95);
  \draw[zust] (2.5,8.25) -- (2.5,9.95);
  % Schutzmechanismen
  \node[ba block hervor, minimum height=6mm, minimum width=108mm, font=\footnotesize] (SM) at (10.1,5.95)
    {Schutzmechanismen: begrenzen \textperiodcentered{} vermeiden \textperiodcentered{} überwachen \textperiodcentered{} anhalten};
  \draw[ba duenn, color=baGruen] (5.9,6.55) -- (5.9,6.25);
  \draw[ba duenn, color=baGruen] (10.55,6.55) -- (10.55,6.25);
  \draw[ba duenn, color=baGruen] (14.35,6.55) -- (14.35,6.25);
  % ---------------- Ebene 3: Herstellerschnittstelle ----------------
  \node[ba block, fill=baHellgrau!40, minimum height=8mm, minimum width=143mm] (DDS) at (8.4,4.2)
    {Herstellerschnittstelle (Unitree-DDS)};
  % Befehle nach unten, Zustände nach oben zwischen Funktionsknoten und Schnittstelle
  \draw[zust] (2.5,4.6) -- (2.5,6.55);
  \draw[ziel] (5.6,5.65) -- (5.6,4.6);
  \draw[zust] (6.2,4.6) -- (6.2,5.65);
  \draw[ziel] (9.5,5.65) -- (9.5,4.6);
  \draw[zust] (10.1,4.6) -- (10.1,5.65);
  \draw[ziel] (11.5,5.65) -- (11.5,4.6);
  \draw[zust] (12.1,4.6) -- (12.1,5.65);
  \draw[ziel] (14.05,5.65) -- (14.05,4.6);
  \draw[zust] (14.65,4.6) -- (14.65,5.65);
  % ---------------- Umschalter ----------------
  \node[draw=baAkzent, ba normal, fill=white, rounded corners=1.5pt, text=black, align=center,
        minimum height=9mm, inner sep=3pt] (SW) at (8.4,2.75) {Umschalter\\Simulation / Realbetrieb};
  \draw[baAkzent, ba normal] (8.4,3.8) -- (SW.north);
  % ---------------- Ebene 4: Simulation oder Roboter ----------------
  \draw[ba gruppe] (1.3,-1.1) rectangle (7.0,1.5);
  \node[klein, anchor=north west] at (1.4,1.5) {Container Simulation};
  \node[ba block, minimum width=0pt, text width=22mm] (SIM) at (2.9,-0.1) {MuJoCo-Simulation};
  \node[ba block, minimum width=0pt, text width=22mm] (SH) at (5.5,-0.1) {Simulierte Hände};
  \draw[ba gruppe] (9.8,-1.1) rectangle (15.5,1.5);
  \node[klein, anchor=north west] at (9.9,1.5) {Realbetrieb};
  \node[ba block, minimum width=0pt, text width=22mm] (RB) at (11.4,-0.1) {Unitree G1 mit Onboard-Regler};
  \node[ba block, minimum width=0pt, text width=22mm] (RH) at (14.0,-0.1) {Inspire-Hände};
  \draw[baAkzent, ba normal, <->] (SW.west) -| (4.15,1.5);
  \draw[baAkzent, ba normal, <->] (SW.east) -| (12.65,1.5);
  \node[font=\footnotesize, anchor=east, align=right] at (4.05,2.1) {Domain 1,\\Schnittstelle Loopback};
  \node[font=\footnotesize, anchor=west, align=left] at (12.75,2.1) {Domain 0,\\Netzwerkkarte};
  % Hände real: eigener Anschluss neben der Herstellerschnittstelle
  \draw[ba normal, dash pattern=on 1pt off 2pt, <->] (HA.east) -- (15.85,7.4) -- (15.85,-0.1) -- (RH.east);
  \node[klein, anchor=east, align=right, fill=white, inner sep=1pt] at (15.8,3.2) {eigener Anschluss,\\nur Realbetrieb};
  % Legende
  \draw[ziel] (1.3,-1.75) -- (2.2,-1.75);
  \node[klein, anchor=west] at (2.3,-1.75) {Ziele und Befehle};
  \draw[zust] (6.0,-1.75) -- (6.9,-1.75);
  \node[klein, anchor=west] at (7.0,-1.75) {Zustände und Meldungen};
\end{tikzpicture}

  \caption[Gesamtarchitektur der Umgebung]{Gesamtarchitektur der Umgebung. Die Bedienwege senden Ziele an die Funktionsknoten, die über die Herstellerschnittstelle wahlweise die Simulation oder den realen Roboter ansprechen.}
  \label{fig:architektur}
\end{figure}

% ---- 6.3.1 Betriebszustände und Übergänge (Kern) ----------------------------------------------
% Text davor: Abbildung~\ref{fig:betriebszustaende}, zusammen mit tab:haltstufen.
\begin{figure}[htbp]
  \centering
  % Betriebszustände und Übergänge der Umgebung (abstrakt, gilt für Simulation und Realbetrieb).
% Eigene Darstellung (Art E, Kap. 6.3.1). Keine Quelle in der Caption. Fasst vier gekoppelte Zustandsmaschinen
% des Codes (Onboard-Regler am Roboter, Ersatzregler der Simulation, Armregler, Bedienmodus) zu einem Diagramm zusammen.
% Basis, je Zustand/Übergang (G1@25d5c1a, Pfade relativ zu g1pilot/):
%   Gedämpft -> Bereit ("Start"): docs/50-A-Echter-Roboter.md F Schritt 2 (Start quittiert und stellt aus Damp auf
%     Preparation); docs/30-T-Stehen-und-Gehen.md Zustandsmaschine (/g1pilot/start(True) -> HOLD, Sim) und
%     Tabelle loco_client (/g1pilot/start).
%   Bereit -> Balancierend (Bedienmodus wählen): docs/30-T-Stehen-und-Gehen.md (start_balancing -> STAND, start_walking
%     -> _enter_walk; real: Tabelle loco_client, Bodenkontakt vor "Walk", knee_load_min_nm); docs/10-A-Demo-GUI.md
%     Abschnitt 1 (genau eine Kachel aktiv); docs/50-A-Echter-Roboter.md E (GREIFEN/LAUFEN nur aus Preparation).
%   GREIFEN <-> GEHEN: docs/30-T-Stehen-und-Gehen.md (START BALANCING im WALK: Ausbremsen, dann STAND; START WALKING).
%   Totmann losgelassen: docs/51-A-Notaus.md Stufen ("Gehen stoppen"); docs/30-T loco_client (Deadman/Timeout).
%   Not-Halt von jedem Zustand, Wirkung: docs/51-T-Notaus.md "Wer was tut" (arm_controller, loco_client, loco_sim,
%     Inspire-Bridge); Beine bleiben stehen (real), Simulation: DAMP (docs/51-A-Notaus.md "In der Simulation").
%   Quittieren: docs/51-A-Notaus.md "Quittieren" (Gamepad: geht nicht, absichtlich; danach Arme schlaff bis GREIFEN).
%   Voll-Notaus -> Gedämpft: docs/51-T-Notaus.md (Damp zuerst, nach Voll-Notaus bleibt Damp), docs/51-A-Notaus.md.
%   Sturz erkannt -> Gedämpft (nur Simulation): docs/30-T-Stehen-und-Gehen.md Zustandsmaschine (_fallen -> DAMP).
% Nicht gezeichnet (nicht für beide Welten belegt oder Detail der Tabelle): Damp per Startleiste aus dem Stehen
%   (nur real), Abfangen und Übergabe WALK->STAND (nur Simulation), zweite Bestätigung vor "Walk" (nur real),
%   Sperre bei Hüfte nicht in der Mitte, Fernbedienung des Herstellers.
% Stil: ba-grafik.tex. Einbinden: % Betriebszustände und Übergänge der Umgebung (abstrakt, gilt für Simulation und Realbetrieb).
% Eigene Darstellung (Art E, Kap. 6.3.1). Keine Quelle in der Caption. Fasst vier gekoppelte Zustandsmaschinen
% des Codes (Onboard-Regler am Roboter, Ersatzregler der Simulation, Armregler, Bedienmodus) zu einem Diagramm zusammen.
% Basis, je Zustand/Übergang (G1@25d5c1a, Pfade relativ zu g1pilot/):
%   Gedämpft -> Bereit ("Start"): docs/50-A-Echter-Roboter.md F Schritt 2 (Start quittiert und stellt aus Damp auf
%     Preparation); docs/30-T-Stehen-und-Gehen.md Zustandsmaschine (/g1pilot/start(True) -> HOLD, Sim) und
%     Tabelle loco_client (/g1pilot/start).
%   Bereit -> Balancierend (Bedienmodus wählen): docs/30-T-Stehen-und-Gehen.md (start_balancing -> STAND, start_walking
%     -> _enter_walk; real: Tabelle loco_client, Bodenkontakt vor "Walk", knee_load_min_nm); docs/10-A-Demo-GUI.md
%     Abschnitt 1 (genau eine Kachel aktiv); docs/50-A-Echter-Roboter.md E (GREIFEN/LAUFEN nur aus Preparation).
%   GREIFEN <-> GEHEN: docs/30-T-Stehen-und-Gehen.md (START BALANCING im WALK: Ausbremsen, dann STAND; START WALKING).
%   Totmann losgelassen: docs/51-A-Notaus.md Stufen ("Gehen stoppen"); docs/30-T loco_client (Deadman/Timeout).
%   Not-Halt von jedem Zustand, Wirkung: docs/51-T-Notaus.md "Wer was tut" (arm_controller, loco_client, loco_sim,
%     Inspire-Bridge); Beine bleiben stehen (real), Simulation: DAMP (docs/51-A-Notaus.md "In der Simulation").
%   Quittieren: docs/51-A-Notaus.md "Quittieren" (Gamepad: geht nicht, absichtlich; danach Arme schlaff bis GREIFEN).
%   Voll-Notaus -> Gedämpft: docs/51-T-Notaus.md (Damp zuerst, nach Voll-Notaus bleibt Damp), docs/51-A-Notaus.md.
%   Sturz erkannt -> Gedämpft (nur Simulation): docs/30-T-Stehen-und-Gehen.md Zustandsmaschine (_fallen -> DAMP).
% Nicht gezeichnet (nicht für beide Welten belegt oder Detail der Tabelle): Damp per Startleiste aus dem Stehen
%   (nur real), Abfangen und Übergabe WALK->STAND (nur Simulation), zweite Bestätigung vor "Walk" (nur real),
%   Sperre bei Hüfte nicht in der Mitte, Fernbedienung des Herstellers.
% Stil: ba-grafik.tex. Einbinden: % Betriebszustände und Übergänge der Umgebung (abstrakt, gilt für Simulation und Realbetrieb).
% Eigene Darstellung (Art E, Kap. 6.3.1). Keine Quelle in der Caption. Fasst vier gekoppelte Zustandsmaschinen
% des Codes (Onboard-Regler am Roboter, Ersatzregler der Simulation, Armregler, Bedienmodus) zu einem Diagramm zusammen.
% Basis, je Zustand/Übergang (G1@25d5c1a, Pfade relativ zu g1pilot/):
%   Gedämpft -> Bereit ("Start"): docs/50-A-Echter-Roboter.md F Schritt 2 (Start quittiert und stellt aus Damp auf
%     Preparation); docs/30-T-Stehen-und-Gehen.md Zustandsmaschine (/g1pilot/start(True) -> HOLD, Sim) und
%     Tabelle loco_client (/g1pilot/start).
%   Bereit -> Balancierend (Bedienmodus wählen): docs/30-T-Stehen-und-Gehen.md (start_balancing -> STAND, start_walking
%     -> _enter_walk; real: Tabelle loco_client, Bodenkontakt vor "Walk", knee_load_min_nm); docs/10-A-Demo-GUI.md
%     Abschnitt 1 (genau eine Kachel aktiv); docs/50-A-Echter-Roboter.md E (GREIFEN/LAUFEN nur aus Preparation).
%   GREIFEN <-> GEHEN: docs/30-T-Stehen-und-Gehen.md (START BALANCING im WALK: Ausbremsen, dann STAND; START WALKING).
%   Totmann losgelassen: docs/51-A-Notaus.md Stufen ("Gehen stoppen"); docs/30-T loco_client (Deadman/Timeout).
%   Not-Halt von jedem Zustand, Wirkung: docs/51-T-Notaus.md "Wer was tut" (arm_controller, loco_client, loco_sim,
%     Inspire-Bridge); Beine bleiben stehen (real), Simulation: DAMP (docs/51-A-Notaus.md "In der Simulation").
%   Quittieren: docs/51-A-Notaus.md "Quittieren" (Gamepad: geht nicht, absichtlich; danach Arme schlaff bis GREIFEN).
%   Voll-Notaus -> Gedämpft: docs/51-T-Notaus.md (Damp zuerst, nach Voll-Notaus bleibt Damp), docs/51-A-Notaus.md.
%   Sturz erkannt -> Gedämpft (nur Simulation): docs/30-T-Stehen-und-Gehen.md Zustandsmaschine (_fallen -> DAMP).
% Nicht gezeichnet (nicht für beide Welten belegt oder Detail der Tabelle): Damp per Startleiste aus dem Stehen
%   (nur real), Abfangen und Übergabe WALK->STAND (nur Simulation), zweite Bestätigung vor "Walk" (nur real),
%   Sperre bei Hüfte nicht in der Mitte, Fernbedienung des Herstellers.
% Stil: ba-grafik.tex. Einbinden: \input{Bilder/tikz/betriebszustaende} in einer figure-Umgebung. Nicht skalieren.
\begin{tikzpicture}[ba]
  \tikzset{
    zu/.style={ba block, align=center, minimum height=12mm},
    halt/.style={ba block, draw=baAkzent, line width=1.2pt, align=center, minimum height=12mm},
    ueb/.style={font=\footnotesize, text=black},
    klein/.style={font=\footnotesize, text=baGrau},
  }
  % Rahmen: Not-Halt und Voll-Notaus gelten von jedem Zustand
  \draw[ba gruppe] (0.3,1.85) rectangle (15.7,8.95);
  \node[klein, anchor=south west] at (0.45,1.95) {von jedem Zustand};
  % Zustände
  \node[zu, text width=30mm] (G) at (2.3,7.9) {Gedämpft\\{\footnotesize alle Gelenke weich}};
  \node[zu, text width=44mm] (B) at (8.2,7.9) {Bereit\\{\footnotesize Gelenke gehalten, noch kein Balancieren}};
  \draw[ba gruppe, draw=baGrau] (4.0,2.3) rectangle (15.3,5.7);
  \node[font=\small, anchor=north west] at (4.15,5.65) {Balancierend (genau ein Bedienmodus aktiv)};
  \node[zu, text width=34mm, ba block hervor] (GR) at (6.1,3.95) {GREIFEN\\{\footnotesize Stehen, Arme frei}};
  \node[zu, text width=40mm, ba block hervor] (GE) at (13.0,3.95) {GEHEN\\{\footnotesize Geschwindigkeitsvorgabe}};
  % Übergänge
  \draw[ba pfeil] (G.east) -- node[ueb, above] {Start} (B.west);
  \draw[ba pfeil] (B.south) -- (8.2,5.7);
  \node[ueb, anchor=west] at (8.4,6.45) {Modus wählen (real: nur mit Bodenkontakt)};
  \draw[ba pfeil] ([yshift=2mm]GR.east) -- node[ueb, above] {GEHEN wählen} ([yshift=2mm]GE.west);
  \draw[ba pfeil] ([yshift=-2mm]GE.west) -- node[ueb, below] {GREIFEN wählen} ([yshift=-2mm]GR.east);
  \node[klein, anchor=south, align=center] at (13.0,2.38) {Totmann losgelassen:\\Gehen stoppt};
  % Sturz (nur Simulation)
  \draw[ba pfeil] (4.0,3.95) -| ([xshift=8mm]G.south);
  \node[ueb, anchor=west, align=left] at (3.2,6.45) {Sturz erkannt\\(nur Simulation)};
  % Not-Halt (Sammelübergang vom Rahmen)
  \node[halt, text width=50mm] (NH) at (6.5,0.1) {Not-Halt (verriegelt)\\{\footnotesize Arme schlaff, Hände halten, Taille gehalten, Beine bleiben stehen}};
  \draw[ba pfeil, baAkzent] (6.5,1.85) -- (NH.north);
  \node[ueb, anchor=west, align=left] at (6.65,1.4) {auf jedem\\Bedienweg};
  % Quittieren
  \draw[ba pfeil] (NH.east) -| (14.5,2.3);
  \node[ueb, anchor=south west, text width=47mm] at (9.6,0.35) {Quittieren (nicht am \mbox{Gamepad}); Arme bleiben schlaff, bis GREIFEN gewählt wird};
  % Voll-Notaus
  \node[halt, text width=24mm] (VN) at (2.1,0.1) {Voll-Notaus\\{\footnotesize alles gedämpft}};
  \draw[ba pfeil, baAkzent] (2.1,1.85) -- (VN.north);
  \node[ueb, anchor=west, text width=30mm] at (2.25,1.4) {Gamepad, Tasten gehalten};
  \draw[ba pfeil] (VN.west) -- ++(-0.2,0) |- (G.west);
\end{tikzpicture}
 in einer figure-Umgebung. Nicht skalieren.
\begin{tikzpicture}[ba]
  \tikzset{
    zu/.style={ba block, align=center, minimum height=12mm},
    halt/.style={ba block, draw=baAkzent, line width=1.2pt, align=center, minimum height=12mm},
    ueb/.style={font=\footnotesize, text=black},
    klein/.style={font=\footnotesize, text=baGrau},
  }
  % Rahmen: Not-Halt und Voll-Notaus gelten von jedem Zustand
  \draw[ba gruppe] (0.3,1.85) rectangle (15.7,8.95);
  \node[klein, anchor=south west] at (0.45,1.95) {von jedem Zustand};
  % Zustände
  \node[zu, text width=30mm] (G) at (2.3,7.9) {Gedämpft\\{\footnotesize alle Gelenke weich}};
  \node[zu, text width=44mm] (B) at (8.2,7.9) {Bereit\\{\footnotesize Gelenke gehalten, noch kein Balancieren}};
  \draw[ba gruppe, draw=baGrau] (4.0,2.3) rectangle (15.3,5.7);
  \node[font=\small, anchor=north west] at (4.15,5.65) {Balancierend (genau ein Bedienmodus aktiv)};
  \node[zu, text width=34mm, ba block hervor] (GR) at (6.1,3.95) {GREIFEN\\{\footnotesize Stehen, Arme frei}};
  \node[zu, text width=40mm, ba block hervor] (GE) at (13.0,3.95) {GEHEN\\{\footnotesize Geschwindigkeitsvorgabe}};
  % Übergänge
  \draw[ba pfeil] (G.east) -- node[ueb, above] {Start} (B.west);
  \draw[ba pfeil] (B.south) -- (8.2,5.7);
  \node[ueb, anchor=west] at (8.4,6.45) {Modus wählen (real: nur mit Bodenkontakt)};
  \draw[ba pfeil] ([yshift=2mm]GR.east) -- node[ueb, above] {GEHEN wählen} ([yshift=2mm]GE.west);
  \draw[ba pfeil] ([yshift=-2mm]GE.west) -- node[ueb, below] {GREIFEN wählen} ([yshift=-2mm]GR.east);
  \node[klein, anchor=south, align=center] at (13.0,2.38) {Totmann losgelassen:\\Gehen stoppt};
  % Sturz (nur Simulation)
  \draw[ba pfeil] (4.0,3.95) -| ([xshift=8mm]G.south);
  \node[ueb, anchor=west, align=left] at (3.2,6.45) {Sturz erkannt\\(nur Simulation)};
  % Not-Halt (Sammelübergang vom Rahmen)
  \node[halt, text width=50mm] (NH) at (6.5,0.1) {Not-Halt (verriegelt)\\{\footnotesize Arme schlaff, Hände halten, Taille gehalten, Beine bleiben stehen}};
  \draw[ba pfeil, baAkzent] (6.5,1.85) -- (NH.north);
  \node[ueb, anchor=west, align=left] at (6.65,1.4) {auf jedem\\Bedienweg};
  % Quittieren
  \draw[ba pfeil] (NH.east) -| (14.5,2.3);
  \node[ueb, anchor=south west, text width=47mm] at (9.6,0.35) {Quittieren (nicht am \mbox{Gamepad}); Arme bleiben schlaff, bis GREIFEN gewählt wird};
  % Voll-Notaus
  \node[halt, text width=24mm] (VN) at (2.1,0.1) {Voll-Notaus\\{\footnotesize alles gedämpft}};
  \draw[ba pfeil, baAkzent] (2.1,1.85) -- (VN.north);
  \node[ueb, anchor=west, text width=30mm] at (2.25,1.4) {Gamepad, Tasten gehalten};
  \draw[ba pfeil] (VN.west) -- ++(-0.2,0) |- (G.west);
\end{tikzpicture}
 in einer figure-Umgebung. Nicht skalieren.
\begin{tikzpicture}[ba]
  \tikzset{
    zu/.style={ba block, align=center, minimum height=12mm},
    halt/.style={ba block, draw=baAkzent, line width=1.2pt, align=center, minimum height=12mm},
    ueb/.style={font=\footnotesize, text=black},
    klein/.style={font=\footnotesize, text=baGrau},
  }
  % Rahmen: Not-Halt und Voll-Notaus gelten von jedem Zustand
  \draw[ba gruppe] (0.3,1.85) rectangle (15.7,8.95);
  \node[klein, anchor=south west] at (0.45,1.95) {von jedem Zustand};
  % Zustände
  \node[zu, text width=30mm] (G) at (2.3,7.9) {Gedämpft\\{\footnotesize alle Gelenke weich}};
  \node[zu, text width=44mm] (B) at (8.2,7.9) {Bereit\\{\footnotesize Gelenke gehalten, noch kein Balancieren}};
  \draw[ba gruppe, draw=baGrau] (4.0,2.3) rectangle (15.3,5.7);
  \node[font=\small, anchor=north west] at (4.15,5.65) {Balancierend (genau ein Bedienmodus aktiv)};
  \node[zu, text width=34mm, ba block hervor] (GR) at (6.1,3.95) {GREIFEN\\{\footnotesize Stehen, Arme frei}};
  \node[zu, text width=40mm, ba block hervor] (GE) at (13.0,3.95) {GEHEN\\{\footnotesize Geschwindigkeitsvorgabe}};
  % Übergänge
  \draw[ba pfeil] (G.east) -- node[ueb, above] {Start} (B.west);
  \draw[ba pfeil] (B.south) -- (8.2,5.7);
  \node[ueb, anchor=west] at (8.4,6.45) {Modus wählen (real: nur mit Bodenkontakt)};
  \draw[ba pfeil] ([yshift=2mm]GR.east) -- node[ueb, above] {GEHEN wählen} ([yshift=2mm]GE.west);
  \draw[ba pfeil] ([yshift=-2mm]GE.west) -- node[ueb, below] {GREIFEN wählen} ([yshift=-2mm]GR.east);
  \node[klein, anchor=south, align=center] at (13.0,2.38) {Totmann losgelassen:\\Gehen stoppt};
  % Sturz (nur Simulation)
  \draw[ba pfeil] (4.0,3.95) -| ([xshift=8mm]G.south);
  \node[ueb, anchor=west, align=left] at (3.2,6.45) {Sturz erkannt\\(nur Simulation)};
  % Not-Halt (Sammelübergang vom Rahmen)
  \node[halt, text width=50mm] (NH) at (6.5,0.1) {Not-Halt (verriegelt)\\{\footnotesize Arme schlaff, Hände halten, Taille gehalten, Beine bleiben stehen}};
  \draw[ba pfeil, baAkzent] (6.5,1.85) -- (NH.north);
  \node[ueb, anchor=west, align=left] at (6.65,1.4) {auf jedem\\Bedienweg};
  % Quittieren
  \draw[ba pfeil] (NH.east) -| (14.5,2.3);
  \node[ueb, anchor=south west, text width=47mm] at (9.6,0.35) {Quittieren (nicht am \mbox{Gamepad}); Arme bleiben schlaff, bis GREIFEN gewählt wird};
  % Voll-Notaus
  \node[halt, text width=24mm] (VN) at (2.1,0.1) {Voll-Notaus\\{\footnotesize alles gedämpft}};
  \draw[ba pfeil, baAkzent] (2.1,1.85) -- (VN.north);
  \node[ueb, anchor=west, text width=30mm] at (2.25,1.4) {Gamepad, Tasten gehalten};
  \draw[ba pfeil] (VN.west) -- ++(-0.2,0) |- (G.west);
\end{tikzpicture}

  \caption[Betriebszustände und Übergänge]{Betriebszustände und Übergänge der Umgebung. Das abstrakte Diagramm gilt für Simulation und Realbetrieb; in der Simulation dämpft der Not-Halt auch die Beine, und der Voll-Notaus wirkt wie der Not-Halt.}
  \label{fig:betriebszustaende}
\end{figure}

% ---- 6.3.3 Übergabe der Arme (Kann) -----------------------------------------------------------
% Text davor: Abbildung~\ref{fig:armuebergabe}, schematisch ohne Messwerte; ein gemessener Verlauf gehört nach Kap. 8.
\begin{figure}[htbp]
  \centering
  % Übernahme und Rückgabe der Arme (Überblendung per Gewicht), schematisch. Keine Messwerte, keine Skalen.
% Eigene Darstellung (Art E, Kap. 6.3.3). Keine Quelle in der Caption.
% Basis: G1@6a40e47 g1pilot/docs/20-T-Arm-Manipulation.md, Abschnitte "rt/arm_sdk-Gewichtsrampe" (Gewicht 0..1,
%   Rampe über arm_weight_ramp_up_s, real 2,0 s, "beim DISABLE symmetrisch zurück") und "Übergabe an den
%   Onboard-Regler" (_record_onboard_pose merkt die gemessene Armpose; DISABLE -> geplante Bahn in die gemerkte
%   Pose, erst bei Ankunft Rampe 1->0; Zeitgrenze release_timeout_s 12 s, hier nicht gezeichnet);
%   docs/50-A-Echter-Roboter.md E ("Übernehmen ohne Bewegung": gemessene Pose = Pose des Onboard-Reglers,
%   kommandiert genau diese, Gewicht 0->1 über 2,0 s).
% a) Übernehmen: Gewicht 0 -> 1, Abstand der kommandierten Pose zur gemerkten Onboard-Pose bleibt 0.
% b) Zurückgeben: Gewicht bleibt 1 während der geplanten Rückfahrt (Abstand sinkt auf 0), danach Rampe 1 -> 0.
% Nicht gezeichnet: Not-Halt (Gewicht bleibt 1, Steifigkeit 0) und Programmende (kurze Rückfahrt, dann Rampe).
% Stil: ba-grafik.tex. Einbinden: % Übernahme und Rückgabe der Arme (Überblendung per Gewicht), schematisch. Keine Messwerte, keine Skalen.
% Eigene Darstellung (Art E, Kap. 6.3.3). Keine Quelle in der Caption.
% Basis: G1@6a40e47 g1pilot/docs/20-T-Arm-Manipulation.md, Abschnitte "rt/arm_sdk-Gewichtsrampe" (Gewicht 0..1,
%   Rampe über arm_weight_ramp_up_s, real 2,0 s, "beim DISABLE symmetrisch zurück") und "Übergabe an den
%   Onboard-Regler" (_record_onboard_pose merkt die gemessene Armpose; DISABLE -> geplante Bahn in die gemerkte
%   Pose, erst bei Ankunft Rampe 1->0; Zeitgrenze release_timeout_s 12 s, hier nicht gezeichnet);
%   docs/50-A-Echter-Roboter.md E ("Übernehmen ohne Bewegung": gemessene Pose = Pose des Onboard-Reglers,
%   kommandiert genau diese, Gewicht 0->1 über 2,0 s).
% a) Übernehmen: Gewicht 0 -> 1, Abstand der kommandierten Pose zur gemerkten Onboard-Pose bleibt 0.
% b) Zurückgeben: Gewicht bleibt 1 während der geplanten Rückfahrt (Abstand sinkt auf 0), danach Rampe 1 -> 0.
% Nicht gezeichnet: Not-Halt (Gewicht bleibt 1, Steifigkeit 0) und Programmende (kurze Rückfahrt, dann Rampe).
% Stil: ba-grafik.tex. Einbinden: % Übernahme und Rückgabe der Arme (Überblendung per Gewicht), schematisch. Keine Messwerte, keine Skalen.
% Eigene Darstellung (Art E, Kap. 6.3.3). Keine Quelle in der Caption.
% Basis: G1@6a40e47 g1pilot/docs/20-T-Arm-Manipulation.md, Abschnitte "rt/arm_sdk-Gewichtsrampe" (Gewicht 0..1,
%   Rampe über arm_weight_ramp_up_s, real 2,0 s, "beim DISABLE symmetrisch zurück") und "Übergabe an den
%   Onboard-Regler" (_record_onboard_pose merkt die gemessene Armpose; DISABLE -> geplante Bahn in die gemerkte
%   Pose, erst bei Ankunft Rampe 1->0; Zeitgrenze release_timeout_s 12 s, hier nicht gezeichnet);
%   docs/50-A-Echter-Roboter.md E ("Übernehmen ohne Bewegung": gemessene Pose = Pose des Onboard-Reglers,
%   kommandiert genau diese, Gewicht 0->1 über 2,0 s).
% a) Übernehmen: Gewicht 0 -> 1, Abstand der kommandierten Pose zur gemerkten Onboard-Pose bleibt 0.
% b) Zurückgeben: Gewicht bleibt 1 während der geplanten Rückfahrt (Abstand sinkt auf 0), danach Rampe 1 -> 0.
% Nicht gezeichnet: Not-Halt (Gewicht bleibt 1, Steifigkeit 0) und Programmende (kurze Rückfahrt, dann Rampe).
% Stil: ba-grafik.tex. Einbinden: \input{Bilder/tikz/armuebergabe} in einer figure-Umgebung. Nicht skalieren.
\begin{tikzpicture}[ba]
  \tikzset{
    achse/.style={ba normal, ->, color=baGrau},
    gew/.style={color=baGruen, line width=1.6pt},
    abst/.style={color=baBlau, line width=1.6pt, dash pattern=on 5pt off 2pt},
    ereig/.style={ba hilfslinie},
    elabel/.style={font=\footnotesize, anchor=south, align=center},
    klein/.style={font=\footnotesize, text=baGrau},
    tick/.style={font=\footnotesize, anchor=east},
  }
  % Layout für 12pt-Grundschrift der Arbeit geprüft (Achsen 5,4 cm, Teilbilder weit genug auseinander, damit
  % die Achsenbeschriftung von b) nicht in a) ragt).
  % ======================== a) Übernehmen ========================
  \begin{scope}[shift={(1.35,0)}]
    % Gewicht
    \draw[achse] (0,4.1) -- (0,6.0);
    \draw[achse] (0,4.1) -- (5.4,4.1);
    \node[tick] at (-0.1,4.1) {0};
    \node[tick] at (-0.1,5.7) {1};
    \draw[ba duenn, color=baGrau] (-0.05,5.7) -- (0.05,5.7);
    \node[font=\footnotesize, rotate=90, anchor=south] at (-0.75,5.05) {Gewicht};
    \draw[gew] (0,4.1) -- (1.0,4.1) -- (3.0,5.7) -- (5.2,5.7);
    \node[klein, anchor=west] at (2.6,4.5) {Rampe über 2\,s};
    % Abstand
    \draw[achse] (0,1.0) -- (0,2.9);
    \draw[achse] (0,1.0) -- (5.4,1.0) node[font=\footnotesize, anchor=north east] {Zeit};
    \node[tick] at (-0.1,1.0) {0};
    \node[font=\footnotesize, rotate=90, anchor=south, align=center] at (-0.75,1.95) {Abstand zur\\Onboard-Pose};
    \draw[abst] (0,1.0) -- (5.2,1.0);
    \node[klein, anchor=west, align=left] at (1.3,1.65) {kommandiert wird die\\gemessene Pose};
    % Ereignis
    \draw[ereig] (1.0,0.85) -- (1.0,6.2);
    \node[elabel] at (1.0,6.25) {Übernehmen};
  \end{scope}
  % Trenner
  \draw[ba trenner] (7.1,0.0) -- (7.1,6.9); % mittig zwischen Achsende a) und Beschriftung b)
  % ======================== b) Zurückgeben ========================
  \begin{scope}[shift={(9.0,0)}]
    % Gewicht
    \draw[achse] (0,4.1) -- (0,6.0);
    \draw[achse] (0,4.1) -- (5.4,4.1);
    \node[tick] at (-0.1,4.1) {0};
    \node[tick] at (-0.1,5.7) {1};
    \draw[ba duenn, color=baGrau] (-0.05,5.7) -- (0.05,5.7);
    \node[font=\footnotesize, rotate=90, anchor=south] at (-0.75,5.05) {Gewicht};
    \draw[gew] (0,5.7) -- (3.3,5.7) -- (5.0,4.1);
    \node[klein, anchor=west] at (4.3,5.3) {Rampe};
    % Abstand
    \draw[achse] (0,1.0) -- (0,2.9);
    \draw[achse] (0,1.0) -- (5.4,1.0) node[font=\footnotesize, anchor=north east] {Zeit};
    \node[tick] at (-0.1,1.0) {0};
    \node[font=\footnotesize, rotate=90, anchor=south, align=center] at (-0.75,1.95) {Abstand zur\\Onboard-Pose};
    \draw[abst] (0,2.6) -- (0.8,2.6) .. controls (1.7,2.6) and (2.2,1.0) .. (3.3,1.0) -- (5.0,1.0);
    \node[klein, anchor=west, align=left] at (3.4,2.05) {geplante,\\langsame\\Rückfahrt};
    % Ereignisse
    \draw[ereig] (0.8,0.85) -- (0.8,6.2);
    \node[elabel] at (0.8,6.25) {Zurückgeben};
    \draw[ereig] (3.3,0.85) -- (3.3,6.2);
    \node[elabel] at (3.3,6.25) {Pose erreicht};
  \end{scope}
  % Teiltitel
  \node[ba teiltitel] at (4.05,0.2) {a) Übernehmen};
  \node[ba teiltitel] at (11.7,0.2) {b) Zurückgeben};
  % Legende
  \draw[gew] (1.35,-0.75) -- (2.15,-0.75);
  \node[klein, anchor=west] at (2.25,-0.75) {Gewicht der eigenen Armregelung};
  \draw[abst] (9.0,-0.75) -- (9.8,-0.75);
  \node[klein, anchor=west] at (9.9,-0.75) {Abstand zur Onboard-Pose};
\end{tikzpicture}
 in einer figure-Umgebung. Nicht skalieren.
\begin{tikzpicture}[ba]
  \tikzset{
    achse/.style={ba normal, ->, color=baGrau},
    gew/.style={color=baGruen, line width=1.6pt},
    abst/.style={color=baBlau, line width=1.6pt, dash pattern=on 5pt off 2pt},
    ereig/.style={ba hilfslinie},
    elabel/.style={font=\footnotesize, anchor=south, align=center},
    klein/.style={font=\footnotesize, text=baGrau},
    tick/.style={font=\footnotesize, anchor=east},
  }
  % Layout für 12pt-Grundschrift der Arbeit geprüft (Achsen 5,4 cm, Teilbilder weit genug auseinander, damit
  % die Achsenbeschriftung von b) nicht in a) ragt).
  % ======================== a) Übernehmen ========================
  \begin{scope}[shift={(1.35,0)}]
    % Gewicht
    \draw[achse] (0,4.1) -- (0,6.0);
    \draw[achse] (0,4.1) -- (5.4,4.1);
    \node[tick] at (-0.1,4.1) {0};
    \node[tick] at (-0.1,5.7) {1};
    \draw[ba duenn, color=baGrau] (-0.05,5.7) -- (0.05,5.7);
    \node[font=\footnotesize, rotate=90, anchor=south] at (-0.75,5.05) {Gewicht};
    \draw[gew] (0,4.1) -- (1.0,4.1) -- (3.0,5.7) -- (5.2,5.7);
    \node[klein, anchor=west] at (2.6,4.5) {Rampe über 2\,s};
    % Abstand
    \draw[achse] (0,1.0) -- (0,2.9);
    \draw[achse] (0,1.0) -- (5.4,1.0) node[font=\footnotesize, anchor=north east] {Zeit};
    \node[tick] at (-0.1,1.0) {0};
    \node[font=\footnotesize, rotate=90, anchor=south, align=center] at (-0.75,1.95) {Abstand zur\\Onboard-Pose};
    \draw[abst] (0,1.0) -- (5.2,1.0);
    \node[klein, anchor=west, align=left] at (1.3,1.65) {kommandiert wird die\\gemessene Pose};
    % Ereignis
    \draw[ereig] (1.0,0.85) -- (1.0,6.2);
    \node[elabel] at (1.0,6.25) {Übernehmen};
  \end{scope}
  % Trenner
  \draw[ba trenner] (7.1,0.0) -- (7.1,6.9); % mittig zwischen Achsende a) und Beschriftung b)
  % ======================== b) Zurückgeben ========================
  \begin{scope}[shift={(9.0,0)}]
    % Gewicht
    \draw[achse] (0,4.1) -- (0,6.0);
    \draw[achse] (0,4.1) -- (5.4,4.1);
    \node[tick] at (-0.1,4.1) {0};
    \node[tick] at (-0.1,5.7) {1};
    \draw[ba duenn, color=baGrau] (-0.05,5.7) -- (0.05,5.7);
    \node[font=\footnotesize, rotate=90, anchor=south] at (-0.75,5.05) {Gewicht};
    \draw[gew] (0,5.7) -- (3.3,5.7) -- (5.0,4.1);
    \node[klein, anchor=west] at (4.3,5.3) {Rampe};
    % Abstand
    \draw[achse] (0,1.0) -- (0,2.9);
    \draw[achse] (0,1.0) -- (5.4,1.0) node[font=\footnotesize, anchor=north east] {Zeit};
    \node[tick] at (-0.1,1.0) {0};
    \node[font=\footnotesize, rotate=90, anchor=south, align=center] at (-0.75,1.95) {Abstand zur\\Onboard-Pose};
    \draw[abst] (0,2.6) -- (0.8,2.6) .. controls (1.7,2.6) and (2.2,1.0) .. (3.3,1.0) -- (5.0,1.0);
    \node[klein, anchor=west, align=left] at (3.4,2.05) {geplante,\\langsame\\Rückfahrt};
    % Ereignisse
    \draw[ereig] (0.8,0.85) -- (0.8,6.2);
    \node[elabel] at (0.8,6.25) {Zurückgeben};
    \draw[ereig] (3.3,0.85) -- (3.3,6.2);
    \node[elabel] at (3.3,6.25) {Pose erreicht};
  \end{scope}
  % Teiltitel
  \node[ba teiltitel] at (4.05,0.2) {a) Übernehmen};
  \node[ba teiltitel] at (11.7,0.2) {b) Zurückgeben};
  % Legende
  \draw[gew] (1.35,-0.75) -- (2.15,-0.75);
  \node[klein, anchor=west] at (2.25,-0.75) {Gewicht der eigenen Armregelung};
  \draw[abst] (9.0,-0.75) -- (9.8,-0.75);
  \node[klein, anchor=west] at (9.9,-0.75) {Abstand zur Onboard-Pose};
\end{tikzpicture}
 in einer figure-Umgebung. Nicht skalieren.
\begin{tikzpicture}[ba]
  \tikzset{
    achse/.style={ba normal, ->, color=baGrau},
    gew/.style={color=baGruen, line width=1.6pt},
    abst/.style={color=baBlau, line width=1.6pt, dash pattern=on 5pt off 2pt},
    ereig/.style={ba hilfslinie},
    elabel/.style={font=\footnotesize, anchor=south, align=center},
    klein/.style={font=\footnotesize, text=baGrau},
    tick/.style={font=\footnotesize, anchor=east},
  }
  % Layout für 12pt-Grundschrift der Arbeit geprüft (Achsen 5,4 cm, Teilbilder weit genug auseinander, damit
  % die Achsenbeschriftung von b) nicht in a) ragt).
  % ======================== a) Übernehmen ========================
  \begin{scope}[shift={(1.35,0)}]
    % Gewicht
    \draw[achse] (0,4.1) -- (0,6.0);
    \draw[achse] (0,4.1) -- (5.4,4.1);
    \node[tick] at (-0.1,4.1) {0};
    \node[tick] at (-0.1,5.7) {1};
    \draw[ba duenn, color=baGrau] (-0.05,5.7) -- (0.05,5.7);
    \node[font=\footnotesize, rotate=90, anchor=south] at (-0.75,5.05) {Gewicht};
    \draw[gew] (0,4.1) -- (1.0,4.1) -- (3.0,5.7) -- (5.2,5.7);
    \node[klein, anchor=west] at (2.6,4.5) {Rampe über 2\,s};
    % Abstand
    \draw[achse] (0,1.0) -- (0,2.9);
    \draw[achse] (0,1.0) -- (5.4,1.0) node[font=\footnotesize, anchor=north east] {Zeit};
    \node[tick] at (-0.1,1.0) {0};
    \node[font=\footnotesize, rotate=90, anchor=south, align=center] at (-0.75,1.95) {Abstand zur\\Onboard-Pose};
    \draw[abst] (0,1.0) -- (5.2,1.0);
    \node[klein, anchor=west, align=left] at (1.3,1.65) {kommandiert wird die\\gemessene Pose};
    % Ereignis
    \draw[ereig] (1.0,0.85) -- (1.0,6.2);
    \node[elabel] at (1.0,6.25) {Übernehmen};
  \end{scope}
  % Trenner
  \draw[ba trenner] (7.1,0.0) -- (7.1,6.9); % mittig zwischen Achsende a) und Beschriftung b)
  % ======================== b) Zurückgeben ========================
  \begin{scope}[shift={(9.0,0)}]
    % Gewicht
    \draw[achse] (0,4.1) -- (0,6.0);
    \draw[achse] (0,4.1) -- (5.4,4.1);
    \node[tick] at (-0.1,4.1) {0};
    \node[tick] at (-0.1,5.7) {1};
    \draw[ba duenn, color=baGrau] (-0.05,5.7) -- (0.05,5.7);
    \node[font=\footnotesize, rotate=90, anchor=south] at (-0.75,5.05) {Gewicht};
    \draw[gew] (0,5.7) -- (3.3,5.7) -- (5.0,4.1);
    \node[klein, anchor=west] at (4.3,5.3) {Rampe};
    % Abstand
    \draw[achse] (0,1.0) -- (0,2.9);
    \draw[achse] (0,1.0) -- (5.4,1.0) node[font=\footnotesize, anchor=north east] {Zeit};
    \node[tick] at (-0.1,1.0) {0};
    \node[font=\footnotesize, rotate=90, anchor=south, align=center] at (-0.75,1.95) {Abstand zur\\Onboard-Pose};
    \draw[abst] (0,2.6) -- (0.8,2.6) .. controls (1.7,2.6) and (2.2,1.0) .. (3.3,1.0) -- (5.0,1.0);
    \node[klein, anchor=west, align=left] at (3.4,2.05) {geplante,\\langsame\\Rückfahrt};
    % Ereignisse
    \draw[ereig] (0.8,0.85) -- (0.8,6.2);
    \node[elabel] at (0.8,6.25) {Zurückgeben};
    \draw[ereig] (3.3,0.85) -- (3.3,6.2);
    \node[elabel] at (3.3,6.25) {Pose erreicht};
  \end{scope}
  % Teiltitel
  \node[ba teiltitel] at (4.05,0.2) {a) Übernehmen};
  \node[ba teiltitel] at (11.7,0.2) {b) Zurückgeben};
  % Legende
  \draw[gew] (1.35,-0.75) -- (2.15,-0.75);
  \node[klein, anchor=west] at (2.25,-0.75) {Gewicht der eigenen Armregelung};
  \draw[abst] (9.0,-0.75) -- (9.8,-0.75);
  \node[klein, anchor=west] at (9.9,-0.75) {Abstand zur Onboard-Pose};
\end{tikzpicture}

  \caption[Übernahme und Rückgabe der Arme]{Schematischer Zeitverlauf der Übernahme und der Rückgabe der Arme am realen Roboter. Das Gewicht blendet zwischen dem Onboard-Regler und der eigenen Armregelung über; die Verläufe sind qualitativ und enthalten keine Messwerte.}
  \label{fig:armuebergabe}
\end{figure}

% ---- 6.4.3 Posen und Abläufe (Kann) -----------------------------------------------------------
% Text davor: Abbildung~\ref{fig:posen_ablaeufe}.
\begin{figure}[htbp]
  \centering
  % Datenmodell der Posen und Abläufe: Pose und relative Bewegung als Bausteine, Ablauf als Folge, Ausführung über die Planung.
% Eigene Darstellung (Art E, Kap. 6.4.3). Keine Quelle in der Caption.
% Basis: G1@25d5c1a g1pilot/docs/23-T-Posen-und-Ablaeufe.md, Abschnitte "Dateiformat" (Pose, relative Bewegung "moves",
%   Ablauf "sequences" verweist nur auf Posennamen und Bewegungen, dieselbe Pose in mehreren Abläufen, wait_s = Pause,
%   Verschiebung je Achse höchstens 0,30 m), "Speicherort" (Datei außerhalb des Containers, bind-gemountet) und
%   "Ablauf-Zustand" (Ausführung am Stück / Schritt für Schritt durch den SequenceRunner der Bedienoberfläche;
%   der Armregler sieht nur einzelne Ziele, jeder Schritt geplant, mit Kollisionsschutz);
%   docs/23-A-Posen-und-Ablaeufe.md (Einlernen durch Hinfahren und Speichern; Schritt für Schritt im Realbetrieb voreingestellt);
%   docs/20-T-Arm-Manipulation.md Abschnitt "Positionsspeicher (plan-execute)" (_start_planned_motion, Planer).
% Nicht gezeichnet: Grenzwert 0,30 m, Statusfolge accepted/executing/reached, Handy- und Gamepad-Aktionen im Detail.
% Stil: ba-grafik.tex. Einbinden: % Datenmodell der Posen und Abläufe: Pose und relative Bewegung als Bausteine, Ablauf als Folge, Ausführung über die Planung.
% Eigene Darstellung (Art E, Kap. 6.4.3). Keine Quelle in der Caption.
% Basis: G1@25d5c1a g1pilot/docs/23-T-Posen-und-Ablaeufe.md, Abschnitte "Dateiformat" (Pose, relative Bewegung "moves",
%   Ablauf "sequences" verweist nur auf Posennamen und Bewegungen, dieselbe Pose in mehreren Abläufen, wait_s = Pause,
%   Verschiebung je Achse höchstens 0,30 m), "Speicherort" (Datei außerhalb des Containers, bind-gemountet) und
%   "Ablauf-Zustand" (Ausführung am Stück / Schritt für Schritt durch den SequenceRunner der Bedienoberfläche;
%   der Armregler sieht nur einzelne Ziele, jeder Schritt geplant, mit Kollisionsschutz);
%   docs/23-A-Posen-und-Ablaeufe.md (Einlernen durch Hinfahren und Speichern; Schritt für Schritt im Realbetrieb voreingestellt);
%   docs/20-T-Arm-Manipulation.md Abschnitt "Positionsspeicher (plan-execute)" (_start_planned_motion, Planer).
% Nicht gezeichnet: Grenzwert 0,30 m, Statusfolge accepted/executing/reached, Handy- und Gamepad-Aktionen im Detail.
% Stil: ba-grafik.tex. Einbinden: % Datenmodell der Posen und Abläufe: Pose und relative Bewegung als Bausteine, Ablauf als Folge, Ausführung über die Planung.
% Eigene Darstellung (Art E, Kap. 6.4.3). Keine Quelle in der Caption.
% Basis: G1@25d5c1a g1pilot/docs/23-T-Posen-und-Ablaeufe.md, Abschnitte "Dateiformat" (Pose, relative Bewegung "moves",
%   Ablauf "sequences" verweist nur auf Posennamen und Bewegungen, dieselbe Pose in mehreren Abläufen, wait_s = Pause,
%   Verschiebung je Achse höchstens 0,30 m), "Speicherort" (Datei außerhalb des Containers, bind-gemountet) und
%   "Ablauf-Zustand" (Ausführung am Stück / Schritt für Schritt durch den SequenceRunner der Bedienoberfläche;
%   der Armregler sieht nur einzelne Ziele, jeder Schritt geplant, mit Kollisionsschutz);
%   docs/23-A-Posen-und-Ablaeufe.md (Einlernen durch Hinfahren und Speichern; Schritt für Schritt im Realbetrieb voreingestellt);
%   docs/20-T-Arm-Manipulation.md Abschnitt "Positionsspeicher (plan-execute)" (_start_planned_motion, Planer).
% Nicht gezeichnet: Grenzwert 0,30 m, Statusfolge accepted/executing/reached, Handy- und Gamepad-Aktionen im Detail.
% Stil: ba-grafik.tex. Einbinden: \input{Bilder/tikz/posen-ablaeufe} in einer figure-Umgebung. Nicht skalieren.
\begin{tikzpicture}[ba]
  \tikzset{
    bst/.style={ba block, align=center, minimum height=17mm, text width=38mm},
    klein/.style={font=\footnotesize, text=baGrau},
    ueb/.style={font=\footnotesize, text=black},
  }
  % Datenebene
  \draw[ba gruppe] (0.2,-0.55) rectangle (10.0,5.5);
  \node[klein, anchor=south west, align=left] at (0.35,-0.5) {Daten: eine Datei außerhalb des Containers,\\in Simulation und Realbetrieb dieselbe};
  \node[bst] (P) at (2.5,4.0) {Pose\\{\footnotesize gespeicherte Stellung von Armen und Händen, einzeln wählbar}};
  \node[bst] (M) at (2.5,1.5) {Relative Bewegung\\{\footnotesize Verschiebung der Hände von der aktuellen Stelle}};
  \node[bst, ba block hervor] (A) at (7.7,2.75) {Ablauf\\{\footnotesize Folge von Posen und relativen Bewegungen, optional mit Pause}};
  \draw[ba pfeil] (P.east) -- ++(0.55,0) |- ([yshift=4mm]A.west);
  \draw[ba pfeil] (M.east) -- ++(0.55,0) |- ([yshift=-4mm]A.west);
  % Einlernen
  \node[ba block, align=center, minimum height=10mm, text width=34mm] (E) at (2.5,6.75) {Einlernen\\{\footnotesize Hinfahren und Speichern}};
  \draw[ba pfeil] (E.south) -- (P.north);
  % Ausführung
  \node[bst, text width=34mm] (X) at (13.2,2.75) {Ausführung\\{\footnotesize am Stück oder Schritt für Schritt}};
  \draw[ba pfeil] (A.east) -- (X.west);
  \node[ba block, align=center, minimum height=10mm, text width=34mm] (BW) at (13.2,5.7) {Bedienwege\\{\footnotesize Start, Weiter}};
  \draw[ba pfeil] (BW.south) -- (X.north);
  % Armregelung
  \node[bst, ba block hervor, text width=40mm] (AR) at (13.0,-0.9) {Armregelung\\{\footnotesize plant jeden Schritt wie ein einzelnes Ziel, mit Kollisionsprüfung}};
  \draw[ba pfeil] (X.south) -- node[ueb, right] {Einzelziele} (X.south |- AR.north);
\end{tikzpicture}
 in einer figure-Umgebung. Nicht skalieren.
\begin{tikzpicture}[ba]
  \tikzset{
    bst/.style={ba block, align=center, minimum height=17mm, text width=38mm},
    klein/.style={font=\footnotesize, text=baGrau},
    ueb/.style={font=\footnotesize, text=black},
  }
  % Datenebene
  \draw[ba gruppe] (0.2,-0.55) rectangle (10.0,5.5);
  \node[klein, anchor=south west, align=left] at (0.35,-0.5) {Daten: eine Datei außerhalb des Containers,\\in Simulation und Realbetrieb dieselbe};
  \node[bst] (P) at (2.5,4.0) {Pose\\{\footnotesize gespeicherte Stellung von Armen und Händen, einzeln wählbar}};
  \node[bst] (M) at (2.5,1.5) {Relative Bewegung\\{\footnotesize Verschiebung der Hände von der aktuellen Stelle}};
  \node[bst, ba block hervor] (A) at (7.7,2.75) {Ablauf\\{\footnotesize Folge von Posen und relativen Bewegungen, optional mit Pause}};
  \draw[ba pfeil] (P.east) -- ++(0.55,0) |- ([yshift=4mm]A.west);
  \draw[ba pfeil] (M.east) -- ++(0.55,0) |- ([yshift=-4mm]A.west);
  % Einlernen
  \node[ba block, align=center, minimum height=10mm, text width=34mm] (E) at (2.5,6.75) {Einlernen\\{\footnotesize Hinfahren und Speichern}};
  \draw[ba pfeil] (E.south) -- (P.north);
  % Ausführung
  \node[bst, text width=34mm] (X) at (13.2,2.75) {Ausführung\\{\footnotesize am Stück oder Schritt für Schritt}};
  \draw[ba pfeil] (A.east) -- (X.west);
  \node[ba block, align=center, minimum height=10mm, text width=34mm] (BW) at (13.2,5.7) {Bedienwege\\{\footnotesize Start, Weiter}};
  \draw[ba pfeil] (BW.south) -- (X.north);
  % Armregelung
  \node[bst, ba block hervor, text width=40mm] (AR) at (13.0,-0.9) {Armregelung\\{\footnotesize plant jeden Schritt wie ein einzelnes Ziel, mit Kollisionsprüfung}};
  \draw[ba pfeil] (X.south) -- node[ueb, right] {Einzelziele} (X.south |- AR.north);
\end{tikzpicture}
 in einer figure-Umgebung. Nicht skalieren.
\begin{tikzpicture}[ba]
  \tikzset{
    bst/.style={ba block, align=center, minimum height=17mm, text width=38mm},
    klein/.style={font=\footnotesize, text=baGrau},
    ueb/.style={font=\footnotesize, text=black},
  }
  % Datenebene
  \draw[ba gruppe] (0.2,-0.55) rectangle (10.0,5.5);
  \node[klein, anchor=south west, align=left] at (0.35,-0.5) {Daten: eine Datei außerhalb des Containers,\\in Simulation und Realbetrieb dieselbe};
  \node[bst] (P) at (2.5,4.0) {Pose\\{\footnotesize gespeicherte Stellung von Armen und Händen, einzeln wählbar}};
  \node[bst] (M) at (2.5,1.5) {Relative Bewegung\\{\footnotesize Verschiebung der Hände von der aktuellen Stelle}};
  \node[bst, ba block hervor] (A) at (7.7,2.75) {Ablauf\\{\footnotesize Folge von Posen und relativen Bewegungen, optional mit Pause}};
  \draw[ba pfeil] (P.east) -- ++(0.55,0) |- ([yshift=4mm]A.west);
  \draw[ba pfeil] (M.east) -- ++(0.55,0) |- ([yshift=-4mm]A.west);
  % Einlernen
  \node[ba block, align=center, minimum height=10mm, text width=34mm] (E) at (2.5,6.75) {Einlernen\\{\footnotesize Hinfahren und Speichern}};
  \draw[ba pfeil] (E.south) -- (P.north);
  % Ausführung
  \node[bst, text width=34mm] (X) at (13.2,2.75) {Ausführung\\{\footnotesize am Stück oder Schritt für Schritt}};
  \draw[ba pfeil] (A.east) -- (X.west);
  \node[ba block, align=center, minimum height=10mm, text width=34mm] (BW) at (13.2,5.7) {Bedienwege\\{\footnotesize Start, Weiter}};
  \draw[ba pfeil] (BW.south) -- (X.north);
  % Armregelung
  \node[bst, ba block hervor, text width=40mm] (AR) at (13.0,-0.9) {Armregelung\\{\footnotesize plant jeden Schritt wie ein einzelnes Ziel, mit Kollisionsprüfung}};
  \draw[ba pfeil] (X.south) -- node[ueb, right] {Einzelziele} (X.south |- AR.north);
\end{tikzpicture}

  \caption[Datenmodell der Posen und Abläufe]{Datenmodell der Posen und Abläufe. Posen und relative Bewegungen sind die Bausteine eines Ablaufs; jeder Schritt der Ausführung durchläuft wie ein einzelnes Ziel die Planung und die Kollisionsprüfung der Armregelung.}
  \label{fig:posen_ablaeufe}
\end{figure}

% ---- 6.5.1 Andockpunkte für externe Projekte (zusätzlich zum Plan, Ergänzung zu tab:andockpunkte) -----
% Text davor: Abbildung~\ref{fig:andockpunkte}; rein konzeptionell, bisher ist kein externes Projekt angebunden (N10).
\begin{figure}[htbp]
  \centering
  % Ebenenmodell der Andockpunkte für externe Projekte: je tiefer ein Projekt andockt, desto mehr Freiheit, desto weniger Schutz.
% Eigene Darstellung (Art E, Kap. 6.5.1), rein konzeptionell; bisher ist kein externes Projekt angebunden (Nils, N10).
% Basis: G1@6a40e47 g1pilot/docs/21-T-Arm-API.md (Architektur: Client -> HTTP -> eigener Knoten -> Armregler, Brücke hält keine
%   Kontrolllogik, Format-Prüfung; Gelenklimits, Kollision und Bereitschaft prüft der Armregler),
%   docs/31-T-Navigation.md (Navigation speist Fahrbefehle über denselben Kanal wie die manuelle Steuerung, joy_mux),
%   docs/02-T-Architektur.md (Zwei Kommunikationsebenen; Sim <-> Real: Ein Schalter; Unitree-DDS identisch zum Roboter),
%   docs/30-T-Stehen-und-Gehen.md und docs/25-T-Eigene-Haende.md (Daten: Lauf-Policies, Hand-Profile nur Simulation),
%   docs/23-T-Posen-und-Ablaeufe.md (Posen und Abläufe als Daten), unitree_mujoco/scene_editor/README.md (Szenen).
% Hinweis: "Schutzmechanismen wirken" bei Arm-API und ROS-2-Topics folgt daraus, dass beide auf demselben Kanal wie die
%   Bedienwege landen (Plan 6.5.1); "wirken nicht" bei der Herstellerschnittstelle direkt: Plan 6.5.1, Befund 3.
% Stil: ba-grafik.tex. Einbinden: % Ebenenmodell der Andockpunkte für externe Projekte: je tiefer ein Projekt andockt, desto mehr Freiheit, desto weniger Schutz.
% Eigene Darstellung (Art E, Kap. 6.5.1), rein konzeptionell; bisher ist kein externes Projekt angebunden (Nils, N10).
% Basis: G1@6a40e47 g1pilot/docs/21-T-Arm-API.md (Architektur: Client -> HTTP -> eigener Knoten -> Armregler, Brücke hält keine
%   Kontrolllogik, Format-Prüfung; Gelenklimits, Kollision und Bereitschaft prüft der Armregler),
%   docs/31-T-Navigation.md (Navigation speist Fahrbefehle über denselben Kanal wie die manuelle Steuerung, joy_mux),
%   docs/02-T-Architektur.md (Zwei Kommunikationsebenen; Sim <-> Real: Ein Schalter; Unitree-DDS identisch zum Roboter),
%   docs/30-T-Stehen-und-Gehen.md und docs/25-T-Eigene-Haende.md (Daten: Lauf-Policies, Hand-Profile nur Simulation),
%   docs/23-T-Posen-und-Ablaeufe.md (Posen und Abläufe als Daten), unitree_mujoco/scene_editor/README.md (Szenen).
% Hinweis: "Schutzmechanismen wirken" bei Arm-API und ROS-2-Topics folgt daraus, dass beide auf demselben Kanal wie die
%   Bedienwege landen (Plan 6.5.1); "wirken nicht" bei der Herstellerschnittstelle direkt: Plan 6.5.1, Befund 3.
% Stil: ba-grafik.tex. Einbinden: % Ebenenmodell der Andockpunkte für externe Projekte: je tiefer ein Projekt andockt, desto mehr Freiheit, desto weniger Schutz.
% Eigene Darstellung (Art E, Kap. 6.5.1), rein konzeptionell; bisher ist kein externes Projekt angebunden (Nils, N10).
% Basis: G1@6a40e47 g1pilot/docs/21-T-Arm-API.md (Architektur: Client -> HTTP -> eigener Knoten -> Armregler, Brücke hält keine
%   Kontrolllogik, Format-Prüfung; Gelenklimits, Kollision und Bereitschaft prüft der Armregler),
%   docs/31-T-Navigation.md (Navigation speist Fahrbefehle über denselben Kanal wie die manuelle Steuerung, joy_mux),
%   docs/02-T-Architektur.md (Zwei Kommunikationsebenen; Sim <-> Real: Ein Schalter; Unitree-DDS identisch zum Roboter),
%   docs/30-T-Stehen-und-Gehen.md und docs/25-T-Eigene-Haende.md (Daten: Lauf-Policies, Hand-Profile nur Simulation),
%   docs/23-T-Posen-und-Ablaeufe.md (Posen und Abläufe als Daten), unitree_mujoco/scene_editor/README.md (Szenen).
% Hinweis: "Schutzmechanismen wirken" bei Arm-API und ROS-2-Topics folgt daraus, dass beide auf demselben Kanal wie die
%   Bedienwege landen (Plan 6.5.1); "wirken nicht" bei der Herstellerschnittstelle direkt: Plan 6.5.1, Befund 3.
% Stil: ba-grafik.tex. Einbinden: \input{Bilder/tikz/andockpunkte-ebenen} in einer figure-Umgebung. Nicht skalieren.
\begin{tikzpicture}[ba]
  \tikzset{
    ebene/.style={ba block, align=left, text width=58mm, minimum width=60mm, minimum height=21mm, anchor=west},
    schutz/.style={ba block hervor, font=\footnotesize, align=center, text width=31mm, minimum width=33mm, minimum height=21mm, anchor=west},
    keinschutz/.style={ba block, font=\footnotesize, draw=baAkzent, line width=1.2pt, align=center, text width=31mm, minimum width=33mm, minimum height=21mm, anchor=west},
    klein/.style={font=\footnotesize, text=baGrau},
  }
  % Raster (Nils, 23.11.2026): Ebenen 1,0..7,0 | Schutz 7,3..10,6 | Daten 10,9..15,3, Abstände je 0,3 cm,
  % Reihen 2,1 cm hoch mit 0,5 cm Abstand; Daten-Kasten bündig mit oberster und unterster Reihe.
  % Pfeil: mehr Freiheit, weniger Schutz
  \draw[baAkzent, line width=1.6pt, ->] (0.3,7.05) -- (0.3,-0.25);
  \node[font=\footnotesize, rotate=90, anchor=south] at (0.35,3.4) {mehr Freiheit \textperiodcentered{} weniger Schutz};
  % Ebene 1: Arm-API
  \node[ebene] (E1) at (1.0,6.0) {\textbf{Arm-API}\\{\footnotesize Zielposen und Gelenkwinkel als Daten über HTTP, kein ROS nötig\par}};
  \node[schutz] (S1) at (7.3,6.0) {Schutzmechanismen wirken};
  % Ebene 2: ROS-2-Topics
  \node[ebene] (E2) at (1.0,3.4) {\textbf{ROS-2-Topics}\\{\footnotesize Ziele über dieselben Kanäle wie die Bedienwege, z.\,B. die Navigation\par}};
  \node[schutz] (S2) at (7.3,3.4) {Schutzmechanismen wirken};
  % Ebene 3: Herstellerschnittstelle direkt
  \node[ebene] (E3) at (1.0,0.8) {\textbf{Herstellerschnittstelle direkt}\\{\footnotesize Projekt spricht den Roboter über das Unitree-SDK an, ohne Änderung auch gegen die Simulation\par}};
  \node[keinschutz] (S3) at (7.3,0.8) {keine Schutz-\\mechanismen\\der Umgebung};
  % Datenebene
  \node[ba block hervor, anchor=north west, text width=40mm, minimum width=44mm, minimum height=73mm, align=left] (D) at (10.9,7.05)
    {\textbf{Daten}\\{\footnotesize unabhängig von der\\Ebene, als Daten\\statt Code:\\[1mm]
     Szenen\\Hand-Profile\\\hspace*{2mm}(nur Simulation)\\Lauf-Policies\\Posen und Abläufe\par}};
\end{tikzpicture}
 in einer figure-Umgebung. Nicht skalieren.
\begin{tikzpicture}[ba]
  \tikzset{
    ebene/.style={ba block, align=left, text width=58mm, minimum width=60mm, minimum height=21mm, anchor=west},
    schutz/.style={ba block hervor, font=\footnotesize, align=center, text width=31mm, minimum width=33mm, minimum height=21mm, anchor=west},
    keinschutz/.style={ba block, font=\footnotesize, draw=baAkzent, line width=1.2pt, align=center, text width=31mm, minimum width=33mm, minimum height=21mm, anchor=west},
    klein/.style={font=\footnotesize, text=baGrau},
  }
  % Raster (Nils, 23.11.2026): Ebenen 1,0..7,0 | Schutz 7,3..10,6 | Daten 10,9..15,3, Abstände je 0,3 cm,
  % Reihen 2,1 cm hoch mit 0,5 cm Abstand; Daten-Kasten bündig mit oberster und unterster Reihe.
  % Pfeil: mehr Freiheit, weniger Schutz
  \draw[baAkzent, line width=1.6pt, ->] (0.3,7.05) -- (0.3,-0.25);
  \node[font=\footnotesize, rotate=90, anchor=south] at (0.35,3.4) {mehr Freiheit \textperiodcentered{} weniger Schutz};
  % Ebene 1: Arm-API
  \node[ebene] (E1) at (1.0,6.0) {\textbf{Arm-API}\\{\footnotesize Zielposen und Gelenkwinkel als Daten über HTTP, kein ROS nötig\par}};
  \node[schutz] (S1) at (7.3,6.0) {Schutzmechanismen wirken};
  % Ebene 2: ROS-2-Topics
  \node[ebene] (E2) at (1.0,3.4) {\textbf{ROS-2-Topics}\\{\footnotesize Ziele über dieselben Kanäle wie die Bedienwege, z.\,B. die Navigation\par}};
  \node[schutz] (S2) at (7.3,3.4) {Schutzmechanismen wirken};
  % Ebene 3: Herstellerschnittstelle direkt
  \node[ebene] (E3) at (1.0,0.8) {\textbf{Herstellerschnittstelle direkt}\\{\footnotesize Projekt spricht den Roboter über das Unitree-SDK an, ohne Änderung auch gegen die Simulation\par}};
  \node[keinschutz] (S3) at (7.3,0.8) {keine Schutz-\\mechanismen\\der Umgebung};
  % Datenebene
  \node[ba block hervor, anchor=north west, text width=40mm, minimum width=44mm, minimum height=73mm, align=left] (D) at (10.9,7.05)
    {\textbf{Daten}\\{\footnotesize unabhängig von der\\Ebene, als Daten\\statt Code:\\[1mm]
     Szenen\\Hand-Profile\\\hspace*{2mm}(nur Simulation)\\Lauf-Policies\\Posen und Abläufe\par}};
\end{tikzpicture}
 in einer figure-Umgebung. Nicht skalieren.
\begin{tikzpicture}[ba]
  \tikzset{
    ebene/.style={ba block, align=left, text width=58mm, minimum width=60mm, minimum height=21mm, anchor=west},
    schutz/.style={ba block hervor, font=\footnotesize, align=center, text width=31mm, minimum width=33mm, minimum height=21mm, anchor=west},
    keinschutz/.style={ba block, font=\footnotesize, draw=baAkzent, line width=1.2pt, align=center, text width=31mm, minimum width=33mm, minimum height=21mm, anchor=west},
    klein/.style={font=\footnotesize, text=baGrau},
  }
  % Raster (Nils, 23.11.2026): Ebenen 1,0..7,0 | Schutz 7,3..10,6 | Daten 10,9..15,3, Abstände je 0,3 cm,
  % Reihen 2,1 cm hoch mit 0,5 cm Abstand; Daten-Kasten bündig mit oberster und unterster Reihe.
  % Pfeil: mehr Freiheit, weniger Schutz
  \draw[baAkzent, line width=1.6pt, ->] (0.3,7.05) -- (0.3,-0.25);
  \node[font=\footnotesize, rotate=90, anchor=south] at (0.35,3.4) {mehr Freiheit \textperiodcentered{} weniger Schutz};
  % Ebene 1: Arm-API
  \node[ebene] (E1) at (1.0,6.0) {\textbf{Arm-API}\\{\footnotesize Zielposen und Gelenkwinkel als Daten über HTTP, kein ROS nötig\par}};
  \node[schutz] (S1) at (7.3,6.0) {Schutzmechanismen wirken};
  % Ebene 2: ROS-2-Topics
  \node[ebene] (E2) at (1.0,3.4) {\textbf{ROS-2-Topics}\\{\footnotesize Ziele über dieselben Kanäle wie die Bedienwege, z.\,B. die Navigation\par}};
  \node[schutz] (S2) at (7.3,3.4) {Schutzmechanismen wirken};
  % Ebene 3: Herstellerschnittstelle direkt
  \node[ebene] (E3) at (1.0,0.8) {\textbf{Herstellerschnittstelle direkt}\\{\footnotesize Projekt spricht den Roboter über das Unitree-SDK an, ohne Änderung auch gegen die Simulation\par}};
  \node[keinschutz] (S3) at (7.3,0.8) {keine Schutz-\\mechanismen\\der Umgebung};
  % Datenebene
  \node[ba block hervor, anchor=north west, text width=40mm, minimum width=44mm, minimum height=73mm, align=left] (D) at (10.9,7.05)
    {\textbf{Daten}\\{\footnotesize unabhängig von der\\Ebene, als Daten\\statt Code:\\[1mm]
     Szenen\\Hand-Profile\\\hspace*{2mm}(nur Simulation)\\Lauf-Policies\\Posen und Abläufe\par}};
\end{tikzpicture}

  \caption[Ebenenmodell der Andockpunkte]{Ebenenmodell der Andockpunkte für externe Projekte. Je tiefer ein Projekt andockt, desto mehr Freiheit hat es und desto weniger Schutzmechanismen der Umgebung wirken; Daten lassen sich unabhängig von der Ebene ergänzen.}
  \label{fig:andockpunkte}
\end{figure}
